\begin{blocksection}
\question \textbf{Which counter changes?} Execute the prompts in order.
Write each displayed value and draw separate attribute tables for the class
and its two instances. If an expression raises an error, write \textbf{Error}
and continue with the next prompt.

\begin{lstlisting}
class Counter:
    step = 1

    def __init__(self, value):
        self.value = value

    def advance(self):
        self.value += self.step
        return self.value

left = Counter(0)
right = Counter(10)
saved = left.advance
\end{lstlisting}
\end{blocksection}

\begin{parts}
\begin{blocksection}
\part
\begin{lstlisting}
>>> Counter.step = 3
>>> left.step = 2
>>> left.step, right.step, Counter.step
\end{lstlisting}
\begin{solution}[.4in]
\lstinline{(2, 3, 3)}. Only \lstinline{left} has an instance attribute named
\lstinline{step}; \lstinline{right} finds the class attribute.
\end{solution}
\end{blocksection}

\begin{blocksection}
\part
\begin{lstlisting}
>>> saved()
\end{lstlisting}
\begin{solution}[.4in]
\begin{lstlisting}
2
\end{lstlisting}
\end{solution}
\begin{lstlisting}
>>> Counter.advance(right)
\end{lstlisting}
\begin{solution}[.4in]
\begin{lstlisting}
13
\end{lstlisting}
The saved bound method supplies the original \lstinline{left} instance as \lstinline{self}. The class function requires an explicit instance.
\end{solution}
\end{blocksection}

\begin{blocksection}
\part
\begin{lstlisting}
>>> left = right
>>> saved()
\end{lstlisting}
\begin{solution}[.4in]
\begin{lstlisting}
4
\end{lstlisting}
\end{solution}
\begin{lstlisting}
>>> left.value, right.value
\end{lstlisting}
\begin{solution}[.4in]
\begin{lstlisting}
(13, 13)
\end{lstlisting}
Reassigning \lstinline{left} does not change the instance held by
\lstinline{saved}. Both names now refer to the second counter, while the bound
method still refers to the first.
\end{solution}
\end{blocksection}

% Out of scope for mentor08: getattr and hasattr.
% \begin{blocksection}
% \part
% \begin{lstlisting}
% >>> hasattr(Counter, "value"), hasattr(right, "value")
% >>> getattr(right, "advance")()
% >>> Counter.value
% \end{lstlisting}
% \begin{solution}[.5in]
% \begin{lstlisting}
% (False, True)
% 16
% Error
% \end{lstlisting}
% \lstinline{value} belongs to instances, not to the class. Looking up
% \lstinline{advance} with \lstinline{getattr} also produces a bound method.
% \end{solution}
% \end{blocksection}
\end{parts}

\begin{questionmeta}
Suggested Time: 8--10 min; Difficulty: Medium
\begin{itemize}
    \item Draw a bound method as a reference to both a function and an instance.
    It retains the instance even if the name used to obtain it is reassigned.
    \item Track instance attributes separately from class attributes. Class
    lookup does not search its instances for their attributes.
\end{itemize}
\end{questionmeta}
